% ===================================================================== PREAMBULE COMMUN — Carnet d'entraînement Fiche 9
\documentclass[11pt,a4paper]{article}
\usepackage[utf8]{inputenc}
\usepackage[T1]{fontenc}
\usepackage{lmodern}
\usepackage[french]{babel}
\usepackage{amsmath,amssymb,mathtools}
\usepackage{icomma}
\usepackage{siunitx}
\sisetup{output-decimal-marker={,},per-mode=symbol}
\DeclareSIUnit{\molar}{M}
\DeclareSIUnit{\Molar}{mol\,L^{-1}}
\usepackage[version=4]{mhchem}
\usepackage{xcolor}
\usepackage{colortbl}
\usepackage{tikz}
\usepackage{pgfplots}
\usepackage[most]{tcolorbox}
\usepackage{enumitem}
\usepackage{array}
\usepackage{booktabs}
\usepackage{multicol}
\usepackage{fancyhdr}
\usepackage[hidelinks]{hyperref}
\usetikzlibrary{arrows.meta,positioning,calc,decorations.pathmorphing,
  decorations.markings,angles,quotes,patterns,backgrounds,
  shapes.geometric,shapes.misc,fadings,intersections,fit}
\pgfplotsset{compat=1.18}
\usepackage[a4paper,margin=2.2cm,headheight=26pt,headsep=10pt]{geometry}

% ---------- Palette ----------
\definecolor{primary}{RGB}{150,98,162}
\definecolor{primarydark}{RGB}{92,52,110}
\definecolor{defcol}{RGB}{0,114,178}
\definecolor{thmcol}{RGB}{0,140,120}
\definecolor{formulacol}{RGB}{198,93,9}
\definecolor{remarkcol}{RGB}{120,94,200}
\definecolor{attncol}{RGB}{200,40,55}
\definecolor{excol}{RGB}{0,135,90}
\definecolor{methcol}{RGB}{90,90,100}
\definecolor{cationcol}{RGB}{0,90,190}
\definecolor{anioncol}{RGB}{200,60,55}
\definecolor{cristalcol}{RGB}{110,105,120}
\definecolor{eaucol}{RGB}{120,180,220}
\definecolor{solcol}{RGB}{0,135,90}
\definecolor{kpscol}{RGB}{198,93,9}
\definecolor{qcol}{RGB}{120,94,200}
\definecolor{precipcol}{RGB}{110,70,40}
\definecolor{exercol}{RGB}{150,98,162}
\definecolor{corrcol}{RGB}{0,120,100}

% ---------- Macros ----------
\newcommand{\Kps}{K_{\mathrm{ps}}}
\newcommand{\Ce}[1]{\left[\,\ce{#1}\,\right]}

% ---------- Style des boîtes ----------
\tcbset{boxbase/.style={enhanced,breakable,boxrule=0.9pt,arc=2.5pt,
   left=3mm,right=3mm,top=2.3mm,bottom=2.3mm,
   fonttitle=\bfseries,coltitle=white,toptitle=0.6mm,bottomtitle=0.6mm}}

\newtcolorbox[auto counter]{definition}[1][]{%
  boxbase,colback=defcol!5,colframe=defcol,
  title={\textsc{Définition}~\thetcbcounter\ \ifx\relax#1\relax\relax\else\textmd{--- \itshape #1}\fi}}
\newtcolorbox[auto counter]{theoreme}[1][]{%
  boxbase,colback=thmcol!6,colframe=thmcol,
  title={\textsc{Loi / Propriété}~\thetcbcounter\ \ifx\relax#1\relax\relax\else\textmd{--- \itshape #1}\fi}}
\newtcolorbox[auto counter]{formule}[1][]{%
  boxbase,colback=formulacol!6,colframe=formulacol,
  title={\textsc{Formule}~\thetcbcounter\ \ifx\relax#1\relax\relax\else\textmd{--- \itshape #1}\fi}}
\newtcolorbox{remarque}[1][]{%
  boxbase,colback=remarkcol!6,colframe=remarkcol,
  title={\textsc{Remarque}\ifx\relax#1\relax\relax\else\textmd{ : \itshape #1}\fi}}
\newtcolorbox{attention}[1][]{%
  boxbase,colback=attncol!6,colframe=attncol,
  title={\textsc{Attention}\ifx\relax#1\relax\relax\else\textmd{ : \itshape #1}\fi}}
\newtcolorbox{exemple}[1][]{%
  boxbase,colback=excol!5,colframe=excol,
  title={\textsc{Exemple}\ifx\relax#1\relax\relax\else\textmd{ : \itshape #1}\fi}}
\newtcolorbox{methode}[1][]{%
  boxbase,colback=methcol!7,colframe=methcol,sharp corners,
  title={\textsc{Méthode}\ifx\relax#1\relax\relax\else\textmd{ : \itshape #1}\fi}}
\newtcolorbox{astuce}[1][]{%
  boxbase,colback=primary!7,colframe=primary,
  title={\textsc{Astuce}\ifx\relax#1\relax\relax\else\textmd{ : \itshape #1}\fi}}

\newtcolorbox[auto counter]{exercice}[1][]{%
  boxbase,colback=exercol!4,colframe=exercol,
  title={\textsc{Exercice}~\thetcbcounter\ \ifx\relax#1\relax\relax\else\textmd{--- \itshape #1}\fi}}
\newtcolorbox[auto counter]{correction}[1][]{%
  boxbase,colback=corrcol!4,colframe=corrcol,
  title={\textsc{Corrigé}~\thetcbcounter\ \ifx\relax#1\relax\relax\else\textmd{--- \itshape #1}\fi}}

% ---------- Environnement « où : » ----------
\newenvironment{ou}{%
  \par\smallskip\noindent\textbf{où :}\nopagebreak
  \begin{itemize}[leftmargin=1.8em,topsep=2pt,itemsep=1.5pt,parsep=0pt]}%
  {\end{itemize}}

% ---------- Styles TikZ ----------
\tikzset{
  cat/.style={circle,fill=cationcol,draw=cationcol!50!black,inner sep=0pt,
              minimum size=5mm,font=\bfseries\scriptsize,text=white},
  an/.style={circle,fill=anioncol,draw=anioncol!50!black,inner sep=0pt,
             minimum size=5mm,font=\bfseries\scriptsize,text=white},
  crx/.style={circle,fill=cristalcol,draw=black!50,inner sep=0pt,
              minimum size=5mm,font=\bfseries\scriptsize,text=white},
  ptn/.style={circle,fill=black,inner sep=1.1pt}}

% ---------- En-têtes ----------
\pagestyle{fancy}
\fancyhf{}
\fancyhead[L]{\small\textcolor{primarydark}{\textbf{Fiche 9} — Solubilité et produit de solubilité}}
\fancyhead[R]{\small\textcolor{primarydark}{\textsc{Carnet d'entraînement}}}
\fancyfoot[C]{\small\thepage}
\renewcommand{\headrulewidth}{0.8pt}
\renewcommand{\headrule}{\hbox to\headwidth{\color{primary}\leaders\hrule height \headrulewidth\hfill}}
\usepackage{titlesec}
\titleformat{\section}{\Large\bfseries\color{primarydark}}{\thesection}{0.6em}{}
\titleformat{\subsection}{\large\bfseries\color{primary}}{\thesubsection}{0.6em}{}

% ---------- Bandeau de titre ----------
% #1 = thème/étiquette   #2 = titre principal   #3 = sous-titre
\newcommand{\bandeau}[3]{%
\begin{center}
\begin{tikzpicture}
\node[rectangle,rounded corners=7pt,fill=primary,text=white,
      minimum width=15.6cm,minimum height=1.95cm,align=center,inner sep=8pt]
  {{\large\bfseries #1}\\[3pt]{\Large\bfseries #2}\\[3pt]{\small #3}};
\end{tikzpicture}
\end{center}
\vspace{2mm}}
\begin{document}
\bandeau{Thème 2 — Relation $\Kps \leftrightarrow s$}%
{Corrigé — Niveau Moyen}%
{Réponses détaillées et justifications}

\begin{correction}[de $\Kps$ à $s$]
\begin{enumerate}[label=\alph*),leftmargin=2em,topsep=2pt,itemsep=2pt]
  \item \textbf{\ce{AgCl}} ($1{:}1$, $\Kps=s^{2}$) :
        $s=\sqrt{\Kps}=\sqrt{\num{1{,}8e-10}}=\boxed{\qty{1{,}3e-5}{\molar}}$.
  \item \textbf{\ce{CaF2}} ($1{:}2$, $\Kps=4\,s^{3}$) :
        $s=\sqrt[3]{\dfrac{\Kps}{4}}=\sqrt[3]{\dfrac{\num{3{,}9e-11}}{4}}
        =\sqrt[3]{\num{9{,}75e-12}}\approx\boxed{\qty{2{,}1e-4}{\molar}}$.
\end{enumerate}
\textbf{Vérification (b)} : $4\,(2{,}1\times10^{-4})^{3}=4\times\num{9{,}3e-12}\approx\num{3{,}7e-11}$,
cohérent avec $\Kps=\num{3{,}9e-11}$.
\end{correction}

\begin{correction}[solubilité massique $\to \Kps$ — \ce{PbBr2}]
$\ce{PbBr2(s) <=> Pb^{2+} + 2 Br-}$.
\begin{enumerate}[label=\alph*),leftmargin=2em,topsep=2pt,itemsep=2pt]
  \item $\Kps=[\ce{Pb^2+}][\ce{Br-}]^{2}=(s)(2s)^{2}=4\,s^{3}$.
  \item Avec $s=s_m/M$ : $\boxed{\ \Kps=4\left(\dfrac{s_m}{M}\right)^{3}
        =4\left(\dfrac{\num{4{,}84}}{M_{\ce{PbBr2}}}\right)^{3}\ }$ \emph{(réponse A du \textsc{qcm}\,n\textsuperscript{o}\,4).}
  \item $s=\dfrac{4{,}84}{367{,}0}=\qty{1{,}32e-2}{\molar}$, donc
        $\Kps=4\,(1{,}32\times10^{-2})^{3}=4\times\num{2{,}30e-6}\approx\boxed{\num{9{,}2e-6}}$.
\end{enumerate}
\end{correction}

\begin{correction}[concentration imposée — \ce{Mn(OH)2}]
$\ce{Mn(OH)2(s) <=> Mn^{2+} + 2 OH-}$.
\begin{enumerate}[label=\alph*),leftmargin=2em,topsep=2pt,itemsep=2pt]
  \item La solubilité est la concentration du cation produit : $s=[\ce{Mn^2+}]=\dfrac{x}{2}$.
  \item $[\ce{OH-}]=2s=2\times\dfrac{x}{2}=x$.
  \item $\Kps=[\ce{Mn^2+}][\ce{OH-}]^{2}=\left(\dfrac{x}{2}\right)(x)^{2}=\boxed{\dfrac{x^{3}}{2}}$.
\end{enumerate}
\emph{Cohérent avec $\Kps=4s^{3}=4(x/2)^{3}=x^{3}/2$ — réponse C du \textsc{qcm}\,n\textsuperscript{o}\,5.}
\end{correction}

\begin{correction}[comparer des sels de même type]
\begin{enumerate}[label=\alph*),leftmargin=2em,topsep=2pt,itemsep=2pt]
  \item Type $3{:}2$ : $[\ce{PO4^3-}]=2s$ et $\Kps=108\,s^{5}$, d'où $s=\sqrt[5]{\Kps/108}$. La fonction
        $\Kps\mapsto s$ est \textbf{croissante}.
  \item Les quatre sels ayant \textbf{la même} relation $\Kps=108\,s^{5}$, on peut comparer directement :
        le plus grand $\Kps$ donne le plus grand $s$, donc le plus grand $[\ce{PO4^3-}]=2s$. C'est
        \boxed{\ce{Sr3(PO4)2}} ($\Kps=\num{1{,}0e-31}$, le plus élevé). \emph{(Réponse C du
        \textsc{qcm}\,n\textsuperscript{o}\,7.)} La comparaison directe n'est légitime que \emph{parce que}
        les sels sont de même type.
\end{enumerate}
\end{correction}

\begin{correction}[comparer des types différents — le piège]
\begin{enumerate}[label=\alph*),leftmargin=2em,topsep=2pt,itemsep=2pt]
  \item \textbf{Non}, il a tort : \ce{AgCl} ($1{:}1$) et \ce{Ag2CrO4} ($2{:}1$) ne suivent pas la même
        loi $\Kps=f(s)$ ($s^2$ contre $4s^3$). On \textbf{ne peut pas} comparer les $\Kps$ directement :
        il faut repasser par $s$.
  \item $s(\ce{AgCl})=\sqrt{\num{1{,}8e-10}}=\qty{1{,}3e-5}{\molar}$ ; \\
        $s(\ce{Ag2CrO4})=\sqrt[3]{\dfrac{\num{1{,}1e-12}}{4}}=\sqrt[3]{\num{2{,}75e-13}}
        =\qty{6{,}5e-5}{\molar}$.\\
        Comme $\qty{6{,}5e-5}{}>\qty{1{,}3e-5}{}$, c'est \boxed{\ce{Ag2CrO4}} qui est le \textbf{plus
        soluble}, alors même que son $\Kps$ est plus petit. Le piège est déjoué.
\end{enumerate}
\end{correction}

\end{document}
